\documentclass[french,11pt,a4paper]{article}

\usepackage[T1]{fontenc}
\usepackage{babel}

\usepackage{geometry}
\geometry{
	a4paper,
	top=2cm,
	bottom=2cm,
	right=3cm,
	left=3cm
}

\usepackage{amsmath}
\usepackage{amsfonts}
\usepackage{amssymb}
\usepackage{amsthm}
\usepackage{nicefrac}
\usepackage{mathrsfs}

\title{Recherches sur la manipulabilité}
\author{hugo}


\newtheorem{definition}{Définition}
\newtheorem{theoreme}{Théorème}


% Ensembles usuels en systèmes de votes
\newcommand{\mc}[1]{\ensuremath{\mathcal{#1}} }  % mc pour MathCal


\begin{document}
	\maketitle
	
	\begin{definition}
		\begin{equation}
			\mathscr{S} : {\mc{D}_\mc{O}}^\mc{V} \rightarrow \mc{O}
		\end{equation}
		Avec $\mc{D}_\mc{O} \in \mc{P}(\mc{O})$. \\
		Remarquons que le système peut être restreint à ${\mc{D}'_\mc{O}}^\mc{V}$ que l'on notera $\mathscr{S}_{|{\mc{D}_\mc{O}}}$. Par exemple, $\mathscr{S}_{|\mathfrak{O}}$. \\
		Un système de vote ou système de choix social est :
		\begin{itemize}
			\item impartial pour les votants si le résultat est invariant par permutation des votants
			\item impartial pour les choix si le résultat est invariant par permutation des choix
			\item  impartial si il est impartial pour les votants et pour les choix
		\end{itemize}
	\end{definition}
	
	\begin{definition}
		Toute relation de préordre induit une relation d'ordre sur l'ensemble des parties choix $\mathcal{P(C)}$. \newline
		Soient $\theta \in \mc{O}$ et $(E,F) \in \mc{P(C)}^2$. Notons $A = \max_\theta E \!\setminus\! (E \cap F)$ et $B = \max_\theta F \!\setminus\! (E \cap F)$ (sous réserve de non-vacuité) (le maximum est un élément quelconque de l'ensemble des majorants).
		\begin{equation}
			E \succ_\theta F \iff
			\left\{ \begin{matrix}
				Vrai &si \ F = \emptyset \\
				Faux &si \ E = \emptyset \ et \ F \neq \emptyset \\
				A \succ_\theta B &si \ A \neq B \\
				E \!\setminus\! \{A\} \succ_\theta F \!\setminus\! \{B\} &si \ A = B
			\end{matrix} \right.
		\end{equation}
	\end{definition}
	
	\begin{proof}
		Cette relation est bien une relation d'ordre/de préordre.
	\end{proof}
	
	\begin{definition}
		Toute relation de préordre induit une relation d'ordre sur l'ensemble des préordres \mc{O}. Soient $\theta \in \mc{O}$ et $(\iota, \kappa) \in \mc{O}^2$
		\begin{equation}
			\iota \succ_\theta \kappa \iff
			\left\{ \begin{matrix}
				Vrai &si \ C = \emptyset \\
				M_\iota(\mc{C}) \succ_\theta M_\kappa(\mc{C}) &si \ M_\iota(\mc{C}) \neq M_\kappa(\mc{C}) \\
				\iota_{|\mc{C} \setminus M_\iota(\mc{C})} \succ_\theta \kappa_{|\mc{C} \setminus M_\kappa(\mc{C})} &si \ M_\iota(\mc{C}) = M_\kappa(\mc{C})
			\end{matrix} \right.
		\end{equation}
	\end{definition}
	
	\begin{definition}
		\begin{equation}
			m_{\mc{C}} \left( \mc{V}, \mc{W} \right) =
			\frac{\left|\left\{ (\theta,\iota) \in \mc{O}^\mc{V} \times \mc{O} \;|\; \exists \iota' \in \mc{O}^\mc{W} : \mathscr{S} (\theta \sqcup \iota') \succnsim_\iota \mathscr{S}(\theta \sqcup \tilde{\iota} \right\}\right|}
			{\left| \mc{O}^\mc{C} \times \mc{O} \right|}
		\end{equation}
	\end{definition}
	
	\begin{definition}
		\begin{equation}
			\forall \left(x,y\right) \in \left(E^F\right)^2\!, \;
			x \sim y
			\iff \exists \sigma \in \mc{S}_F : x = \left(y_{\sigma(f)}\right)_{f \in F}
		\end{equation}
	\end{definition}
	
	\begin{definition}
		\begin{equation}
			\forall \left(\pi,\pi'\right) \in \left(\mc{O}^\mc{V}\right)^2\!, \;
			\pi \sim \pi'
			\iff \exists \sigma \in \mc{S}_\mc{V} : \pi = \left(\pi'_{\sigma(v)}\right)_{v \in \mc{V}}
		\end{equation}
	\end{definition}
	
	\begin{proof}
		$\sim$ est une relation d'équivalence. \\
		Notons $S$ un système de représentant de $\nicefrac{{\mc{D}_\mc{O}}^\mc{V}}{\sim}$. En particulier, $\nicefrac{\mathfrak{O}^\mc{V}}{\sim} = \{\{\mathfrak{O}\}\}$. \\
		\mc{O} peut être vu comme une partie de $\mc{P}(\mc{C}^2)$ (et alors ???). Notons $T$ un système de représentant de $\nicefrac{\mc{O}}{\sim}$.
	\end{proof}
	
	\begin{definition}
		\begin{equation}
			\forall (\theta, \theta') \in \mc{O}^2, \;
			\theta \sim \theta' \iff
			\exists \sigma \in \mc{S}_\mc{V} : \forall (A, B) \in \mc{C}^2,
			\left[ A \succ_\theta B \!\iff\! \sigma(A) \succ_{\theta'} \sigma(B) \right]
		\end{equation}
	\end{definition}
	
	\begin{theoreme}
		\underline{invariance par impartialité} \\
		Soit $\mathscr{S}$ un système de vote impartial.
		\begin{equation}
			m_{\mc{C}} \left( \mc{V}, \mc{W} \right) =
			\frac{1}{\left|\mc{O}\right|^{\left|\mc{V}\right|+1}}
			\sum_{\iota \in S} \left|\bar{\iota}\right|
			\left| \left\{
				\theta \in \mc{O}^\mc{V} \;|\;
				\exists \iota' \in \mc{O}^\mc{W} : \mathscr{S} (\theta \sqcup \iota') \succnsim_\iota \mathscr{S}(\theta \sqcup \tilde{\iota})
			\right\} \right|
		\end{equation}
	\end{theoreme}
	
	\begin{definition} \textit{très générale}
		\begin{equation}
			m (\mc{V}, E, \mc{R}) =
			\frac{
				\left|\left\{
					(\theta, \iota) \in \mc{O}^\mc{V} \times E \;|\;
					\theta \;\mc{R}\; \iota
				\right\}\right|
			}{
				\left| \mc{O}^\mc{V} \times E \right|
			}
		\end{equation}
	\end{definition}
	
	\begin{theoreme}
		\underline{invariance par impartialité} \\
		Soit $\mathscr{S}$ un système de vote impartial.
		\begin{equation}
			m (\mc{V}, E, \mc{R}) =
			\frac{1}{ \left|\mc{O}\right|^{\left|\mc{V}\right|} \left|E\right| }
			\sum_{\iota \in T} \left|\bar{\iota}\right|
			\left(
				\sum_{\theta \in S} 
				\left\{\begin{matrix}
					\left|\bar{\theta}\right| &\text{ si } \theta \;\mc{R}\; \iota \\
					0 &\text{ sinon}
				\end{matrix}\right.
			\right)
		\end{equation}
	\end{theoreme}
	
	\begin{theoreme}
		\begin{equation}
			m_\mathfrak{O}(\mc{V}, E, \mc{R}) = 
			\frac{1}{ 2^{|\mc{C}| \times |\mc{V}|} \cdot |E| }
			\sum_{\iota \in T} \left|\bar{\iota}\right|
			\left(
			\sum_{\theta \in S} 
			\left\{\begin{matrix}
				\left|\bar{\theta}\right| &\text{ si } \theta \;\mc{R}\; \iota \\
				0 &\text{ sinon}
			\end{matrix}\right.
			\right)
		\end{equation}
	\end{theoreme}
	
	\begin{definition}
		\begin{equation}
			\forall \left(A,B\right) \in \mc{C}^2\!, \;
			A \asymp B
			\iff A \succ_\theta B \wedge B \succ_\theta A
		\end{equation}
		$\asymp$ est une relation d'équivalence. Notons $S'_\theta$ un système de représentant.
	\end{definition}
	
	\begin{theoreme}
		\begin{equation}
			\left| \mc{O} \right|
			= \sum_{\iota \in S} \left|\bar{\iota}\right|
				= \sum_{\iota \in S} \frac{
					\left|\mc{C}\right| !
				}{ \displaystyle
					\prod_{A \in S'_\iota} \left|\bar{A}\right| !
				}
		\end{equation}
	\end{theoreme}
	
	\begin{definition}
		\underline{Critère d'honnêteté}
		\begin{equation}
			\forall \iota' \in \mc{O}^\mc{W} : \mathscr{S}(\theta \sqcup \tilde{\iota}) \succ_\iota \mathscr{S} (\theta \sqcup \iota')
		\end{equation}
	\end{definition}
\end{document}